\subsection{Piece (namespace pieces)}
\paragraph{Descrizione modulo:}
Il namespace \texttt{pieces} fornisce funzioni e tabelle pre-calcolate per la generazione
veloce di attacchi e mosse per tutti i tipi di pezzi. Questo modulo è fondamentale per
l'efficienza della generazione mosse e del calcolo di case attaccate.

\paragraph{Tipologie di pezzi:}
Il modulo gestisce tutti e sei i tipi di pezzi:
\begin{itemize}
  \item \textbf{Pedone (Pawn)}: Movimento asimmetrico, dipendente dal colore
  \item \textbf{Cavallo (Knight)}: Movimento a "L", salta altri pezzi
  \item \textbf{Alfiere (Bishop)}: Movimento diagonale, sliding piece
  \item \textbf{Torre (Rook)}: Movimento ortogonale, sliding piece
  \item \textbf{Regina (Queen)}: Combinazione di alfiere e torre
  \item \textbf{Re (King)}: Una casa in tutte le direzioni
\end{itemize}

\paragraph{Funzioni principali:}

\subparagraph{Pedoni:}
\begin{lstlisting}[language=C++]
U64 getPawnAttacks(int16_t squareIndex, bool isWhite);
U64 getPawnForwardPushes(int16_t squareIndex, bool isWhite,
                         U64 occupancy);
U64 getPawnAttackersTo(int16_t targetIndex, bool isWhite);
\end{lstlisting}

\begin{itemize}
  \item \texttt{getPawnAttacks}: Restituisce bitboard delle case attaccate in diagonale
        da un pedone in una data posizione.
  \item \texttt{getPawnForwardPushes}: Restituisce bitboard delle case raggiungibili
        avanzando (1 o 2 case dalla posizione iniziale), considerando l'occupancy.
  \item \texttt{getPawnAttackersTo}: Data una casa target, restituisce bitboard delle
        case da cui un pedone di un certo colore potrebbe attaccare il target (funzione inversa).
\end{itemize}

\subparagraph{Cavalli e Re:}
\begin{lstlisting}[language=C++]
U64 getKnightAttacks(int16_t squareIndex);
U64 getKingAttacks(int16_t squareIndex);
\end{lstlisting}

Restituiscono bitboard delle case attaccabili da cavallo/re da una data posizione.
Non dipendono dall'occupancy poiché il cavallo salta e il re muove di una sola casa.

\subparagraph{Pezzi sliding (Alfiere, Torre, Regina):}
\begin{lstlisting}[language=C++]
U64 getBishopAttacks(int16_t squareIndex, U64 occupancy);
U64 getRookAttacks(int16_t squareIndex, U64 occupancy);
U64 getQueenAttacks(int16_t squareIndex, U64 occupancy);
\end{lstlisting}

Restituiscono bitboard delle case attaccabili considerando l'occupancy della scacchiera.
I raggi si fermano al primo pezzo incontrato (che può essere catturato).

Nota: Queste funzioni usano internamente \textbf{magic bitboards} per performance ottimali.

\paragraph{Lookup Tables pre-calcolate:}

Per massimizzare le performance, il modulo pre-calcola a compile-time tabelle di lookup
per pezzi non-sliding:

\begin{lstlisting}[language=C++]
// Pedoni: [colore][casa] (0=nero, 1=bianco)
constexpr std::array<std::array<uint64_t, 64>, 2> PAWN_ATTACKS;
constexpr std::array<std::array<uint64_t, 64>, 2> PAWN_ATTACKERS_TO;

// Cavalli: [casa]
constexpr std::array<uint64_t, 64> KNIGHT_ATTACKS;

// Re: [casa]
constexpr std::array<uint64_t, 64> KING_ATTACKS;
\end{lstlisting}

Queste tabelle sono inizializzate usando \texttt{constexpr} lambda functions,
quindi il calcolo avviene a compile-time e il risultato è hard-coded nel binario.

\textbf{Vantaggio}: Accesso O(1) a runtime senza calcoli, solo un array lookup.

Esempio di utilizzo:
\begin{lstlisting}[language=C++]
// Ottieni attacchi del cavallo in e4 (indice 28)
uint64_t attacks = pieces::KNIGHT_ATTACKS[28];

// Ottieni attacchi del pedone bianco in e2 (indice 12)
uint64_t pawnAtk = pieces::PAWN_ATTACKS[1][12];  // 1 = white
\end{lstlisting}

\paragraph{Funzioni di utilità:}

\begin{lstlisting}[language=C++]
constexpr int16_t fileOf(int16_t sq);  // Colonna (0-7)
constexpr int16_t rankOf(int16_t sq);  // Riga (0-7)
\end{lstlisting}

Estraggono colonna e riga da un indice di casa (0-63).

\paragraph{Offset per pezzi non-sliding:}

\begin{lstlisting}[language=C++]
constexpr int16_t KNIGHT_OFFSET[8][2] = {
    {1,2}, {2,1}, {2,-1}, {1,-2},
    {-1,-2}, {-2,-1}, {-2,1}, {-1,2}
};

constexpr int16_t KING_OFFSET[8][2] = {
    {1,1}, {1,0}, {1,-1}, {0,-1},
    {-1,-1}, {-1,0}, {-1,1}, {0,1}
};
\end{lstlisting}

Array di offset $(dx, dy)$ per generare le posizioni attaccabili da cavalli e re.

\paragraph{Magic Bitboards (sliding pieces):}

Per sliding pieces (alfiere, torre, regina), il modulo usa la tecnica dei
\textbf{magic bitboards}, che permette di calcolare attacchi in tempo O(1) usando
lookup tables.

\subparagraph{Principio:}
\begin{enumerate}
  \item Per ogni casa e tipo di pezzo (rook/bishop), pre-calcola tutte le possibili
        configurazioni di occupancy sui raggi.
  \item Usa una funzione hash "magica" per mappare configurazioni a indici di tabella.
  \item Memorizza il bitboard di attacchi risultante nella tabella.
  \item A runtime: hasha l'occupancy corrente e accedi direttamente alla tabella.
\end{enumerate}

\textbf{Complessità}:
\begin{itemize}
  \item Spazio: $\sim$800 KB di lookup tables
  \item Tempo: O(1) a runtime (1 hash + 1 array access)
\end{itemize}

\paragraph{Integrazione con Board:}

Il modulo \texttt{pieces} è usato intensivamente da \texttt{Board} per:
\begin{itemize}
  \item \textbf{Generazione mosse pseudo-legali}: \texttt{Board::getPieceMovesBB()}
        chiama le funzioni di attack per ottenere destinazioni possibili.
  \item \textbf{Verifica case attaccate}: \texttt{Board::isSquareAttacked()} usa
        le funzioni per determinare se una casa è sotto attacco.
  \item \textbf{Check detection}: Verifica se il re è attaccato.
  \item \textbf{Pin calculation}: Usa raggi di sliding pieces per trovare pezzi inchiodati.
\end{itemize}

\paragraph{Decisioni progettuali:}

\subparagraph{Namespace invece di classe:}
\texttt{pieces} è un namespace e non una classe perché fornisce solo funzioni pure
e costanti globali, senza bisogno di stato o istanze.

\subparagraph{Constexpr e compile-time computation:}
Le tabelle di lookup sono \texttt{constexpr}, quindi:
\begin{itemize}
  \item Calcolo fatto dal compilatore, non a runtime
  \item Nessun overhead di inizializzazione all'avvio del programma
  \item Tabelle memorizzate in sezione read-only del binario
\end{itemize}

\subparagraph{Bitboard (uint64\_t) come tipo di ritorno:}
Usare bitboard invece di vector permette:
\begin{itemize}
  \item Operazioni batch veloci (AND, OR, XOR bit-per-bit)
  \item Popcount hardware per contare case
  \item Iterazione veloce con operazioni bit-scan
\end{itemize}

\subparagraph{Inline functions:}
Molte funzioni sono dichiarate \texttt{inline} per suggerire al compilatore di
espanderle in-place, eliminando l'overhead di chiamata per funzioni piccole e
frequentemente usate.

\subparagraph{Template di Lambda per array initialization:}
La sintassi:
\begin{lstlisting}[language=C++]
constexpr std::array<uint64_t, 64> TABLE = []{
    std::array<uint64_t, 64> table{};
    // ... calcoli ...
    return table;
}();
\end{lstlisting}

È un idioma C++ moderno che permette inizializzazione complessa di array constexpr
eseguendo una lambda immediatamente (\texttt{()}). Tutto avviene a compile-time.

\paragraph{Performance:}

Grazie a lookup tables e magic bitboards, il modulo \texttt{pieces} è estremamente efficiente:
\begin{itemize}
  \item Knight/King attacks: ~1-2 ns (single array access)
  \item Pawn attacks: ~1-2 ns (single array access)
  \item Bishop/Rook/Queen attacks: ~5-10 ns (hash + array access)
\end{itemize}

Questo è critico poiché queste operazioni sono eseguite milioni di volte durante la ricerca.
